\chapter{Decoupling the Yield Function from the Return Algorithm}

\section{Motivation}

Making the \texttt{YieldFunction} explicit was an important semantic step, but
it did not yet remove the deepest coupling in the local plasticity kernel:
\texttt{PlasticityRelation} still contained a hard-coded radial-return
algorithm.  That coupling is acceptable for classical small-strain \(J_2\)
plasticity, but it becomes a real architectural limitation as soon as the
library needs to support:

\begin{itemize}
  \item cutting-plane or closest-point schemes;
  \item viscoplastic overstress integration;
  \item substepping or adaptive local solves;
  \item pressure-sensitive surfaces with different local correction logic;
  \item damage-plasticity or capped plasticity where the local admissibility
        problem is no longer the same scalar radial-return update.
\end{itemize}

If the return algorithm stays embedded in the constitutive law, then every new
local scheme forces a new constitutive class.  That is exactly the kind of
semantic drift that the previous chapters were trying to remove.

\section{New Separation}

The plasticity architecture now exposes six independent compile-time axes:

\begin{enumerate}
  \item constitutive space / measure policy;
  \item yield criterion;
  \item hardening law;
  \item flow rule;
  \item yield function;
  \item return algorithm.
\end{enumerate}

The new policy lives in \texttt{ReturnAlgorithm.hh} under the concept
\texttt{ReturnAlgorithmPolicy<A, Relation>}.  The contract is intentionally
narrow:

\begin{lstlisting}[language=C++]
algorithm.integrate(relation, total_strain, alpha)
    -> Relation::ReturnMapResultT
\end{lstlisting}

This means the algorithm works on:

\begin{itemize}
  \item a concrete constitutive relation that exposes a typed plasticity
        kernel;
  \item a total kinematic state;
  \item an explicit algorithmic state \(\alpha\).
\end{itemize}

\section{What Remains Inside \texttt{PlasticityRelation}}

\texttt{PlasticityRelation} is still the owner of the constitutive physics:

\begin{itemize}
  \item elastic predictor;
  \item effective trial state for kinematic hardening;
  \item scalar yield-function evaluation;
  \item flow direction;
  \item hardening modulus;
  \item plastic internal-variable evolution;
  \item consistent tangent formula.
\end{itemize}

What moved out is the \emph{algorithmic orchestration} of those pieces.

That is the right cut.  The law keeps the constitutive mathematics; the return
algorithm decides how the local admissibility problem is solved.

\section{Kernel API for Algorithms}

To avoid friend-based hacks or dynamic polymorphism, the relation now exposes a
small public kernel API:

\begin{itemize}
  \item \texttt{elastic\_predictor(...)}
  \item \texttt{effective\_trial(...)}
  \item \texttt{evaluate\_yield\_function(...)}
  \item \texttt{hardening\_modulus(...)}
  \item \texttt{flow\_direction(...)}
  \item \texttt{consistency\_increment(...)}
  \item \texttt{corrected\_stress(...)}
  \item \texttt{evolve\_internal\_variables(...)}
  \item \texttt{consistent\_tangent(...)}
  \item \texttt{make\_elastic\_result(...)}
\end{itemize}

This interface is not intended as user-facing high-level API.  It is an
algorithm-support surface that allows compile-time return algorithms to be
implemented outside the constitutive law without paying runtime dispatch.

\section{Default Algorithm}

The legacy behaviour is preserved by
\texttt{StandardRadialReturnAlgorithm}.  It implements the same backward-Euler
radial-return correction that used to live inside
\texttt{PlasticityRelation}.  Therefore the existing \(J_2\) path remains
backward compatible.

The important architectural change is not a different numerical answer, but the
fact that the local solve is no longer hard-wired to the law itself.

\section{What the New Test Demonstrates}

The regression suite adds a deliberately simple alternative algorithm:
\texttt{HalfStepReturnAlgorithm}.  It reuses the same constitutive kernel but
applies only half of the classical correction increment.  This is not a
production algorithm; it is a proof of separation.

The test checks that:

\begin{itemize}
  \item the custom algorithm satisfies \texttt{ReturnAlgorithmPolicy};
  \item the stress response changes;
  \item the algorithmic tangent changes;
  \item the rest of the constitutive composition remains the same.
\end{itemize}

This is the right kind of regression: it proves that changing the local
algorithm does not require redefining the yield criterion, hardening, flow, or
yield function.

\section{Why This Matters for Future Material Models}

This split is strategically important for the material roadmap discussed in the
material-module analysis chapter.  Reinforced concrete across scales will
eventually need models such as:

\begin{itemize}
  \item concrete damage-plasticity;
  \item fracture-regularized softening;
  \item cyclic degradation with nested loading surfaces;
  \item rate effects and dynamic overstress;
  \item large-strain constitutive updates.
\end{itemize}

Those models are not only ``different yield functions''.  They often require
different \emph{local integration algorithms}.  By making that axis explicit
now, the architecture stops conflating:

\begin{itemize}
  \item the surface being checked;
  \item the scalar admissibility map;
  \item the nonlinear method used to restore admissibility.
\end{itemize}

\section{Immediate Boundary}

The current split is honest but still partial:

\begin{itemize}
  \item the default tangent formula remains the standard small-strain
        rate-independent \(J_2\) formula;
  \item the kernel methods exposed by \texttt{PlasticityRelation} still encode
        assumptions specific to that family;
  \item more general return algorithms may need richer kernel contracts or even
        different state objects.
\end{itemize}

That is acceptable.  The present goal is not to solve every future material
model today; it is to stop blocking them with a hard-coded radial return.

\section{Recommended Next Step}

The next correct move is to apply the same discipline to richer local models:

\begin{itemize}
  \item introduce a non-associative or pressure-sensitive return algorithm;
  \item let the local algorithm request its own residual/Jacobian assembly if
        the scalar closed form is not available;
  \item eventually separate \texttt{ConstitutiveKernel} from
        \texttt{ConstitutiveIntegrator} for continuum damage/fracture models.
\end{itemize}

At that point the library will have a genuinely extensible foundation for
developing nonlinear concrete models beyond classical metal plasticity.
